\section{Introduction}
LLM-powered agentic systems~\citep{autogpt,langgraph,wu2024autogen}, which integrate LLMs' reasoning capabilities with external functionalities, have exhibited impressive performance and facilitated a surge of practical applications like information search~\citep{li2025search}, data analysis~\citep{TimeART}, code generation~\citep{zheng2023codegeex}, deep research~\citep{zheng2025deepresearcher} and autonomous driving~\citep{jin2023surrealdriver}. To further improve the performance and robustness in complex scenarios, multi-agent systems (MAS) are studied to achieve the collaborative intelligence. Generally, multi-agent systems~\citep{ding2024large,han2024llm} are customized for a certain type of task, with a meticulously-designed workflow or communication topology, focusing on how agents effectively utilize, exchange, and induce information reciprocally.
\begin{figure}[t]
    \centering
\includegraphics[width=1\linewidth]{Figures/intro.pdf}
    \caption{Existing multi-agent systems. Our proposed \ourmethod pioneers the exploration of Spatio-Temporal Evolving multi-agent systems.}
\label{fig: intro}
\end{figure}

Classic multi-agent systems often adopt static topologies--see Figure~\ref{fig: intro} (top left), which organizes agents as sequential workflows (chain)~\citep{wei2022chain,hong2023metagpt}, task splitting mode (tree)~\citep{yao2023tree,yao2022react}, or task assignment mode (star)~\citep{wu2024autogen}. Since they are all special cases of graphs, some recent multi-agent systems~\citep{langgraph,wu2024autogen} adopt graphs as their template structures and demonstrate remarkable performance. However, these task-specific static designs rely on expert intervation, and lack automation and adaptation when scenario changes. Therefore, current research interests gradually shift towards dynamic workflow orchestration~\citep{niu2025flow,zhangaflow,lee2025confidence} and communication topology design~\citep{GDesigner,li2025assemble,zhang2025safesieve}, aiming at task-specific or query-specific automated fine-grained customization, known as ``\textit{self-evolving}'' MAS~\citep{fang2025comprehensive,wang2025evoagentx}.

As shown in Figure~\ref{fig: intro}, we categorize existing self-evolving multi-agent systems into ``Spatial Evolving'' and ``Temporal Evolving''. Considering the Spatial Evolving MAS, the mainstream approach is to first formulate the communication topology as a graph, then learn from some cases to train a ``designer'' (often compact neural networks), which can design topologies for specific tasks~\citep{liu2023dynamic,zhuge2024gptswarm} or queries~\citep{GDesigner,li2025assemble,zhang2025safesieve}. However, the designed topologies are fixed throughout the entire conversation process for a query. That means, if a complex query needs multiple dialogue iterations to solve, the communication topology cannot be dynamically tuned during iterations to allocate computing resources and functions among different agent nodes, \textit{which leads to low flexibility and robustness, and may cause high token consumption.} On the other hand, though the Temporal Evolving MAS makes temporal split on workflows, most of them focus on routing single agent in the proper position~\citep{niu2025flow,zhangaflow,lee2025confidence}, \textit{which neglects the potential of collaboration in key temporal iterations.} 

The questions in single-dimension evolution posed above can be explicitly summarized as: \ding{182} \textbf{Performance.} How to fully unleash the collaborative capabilities of MAS in different dialogue iterations? \ding{183} \textbf{Flexibility.} How to make elastic temporal scheduling to fulfill the changing demand? \ding{184} \textbf{Efficiency.} How to ensure efficiency considering the token consumption of MAS and the cost of scheduling? To tackle these problems, we propose \ourmethod, the first MAS with the capability of Spatio-Temporal Evolution, achieving dialogue-wise temporal scheduling for spatial communication topologies. 

The motivation of \ourmethod is to precisely schedule multi-agent collaborations from both temporal and spatial dimensions, preserving performance, flexibility, and efficiency. To achieve this, we believe the core elements lie in: \ding{168} \textbf{Appropriate data structures.} The structure should be able to well maintain the Spatio-Temporal trajectories, providing the basis for scheduling. \ding{171} \textbf{A compact but effective scheduler.} The scheduling is frequent, and the demand for accuracy is high, thus calling for a lightweight and powerful scheduler. \ding{169} \textbf{System state perception.} The scheduler should consider the system state, such as whether it is reliable or not, thus leading to correct scheduling scheme. \ding{170} \textbf{Self-feedback ability.} The system needs to accumulate experience and learn underlying scheduling principles from it.

Based on above-mentioned design principles, \ourmethod adopts \textit{Spatio-Temporal Graph} to maintain the scheduling trajectories, which is an important structure in Spatio-Temporal data mining~\citep{wang2020deep,yang2024survey}, able to characterize the spatial correlations (agent collaboration) at each temporal index (dialogue iteration). Considering the scheduler, we meticulously devise a compact neural network, driven by an MLP-based \textit{Flow-Matching}~\citep{FlowMatch} module, which possesses strong generative mechanism while keeping lightweight, and can work as a stochastic interpolant~\citep{wu2026aurora} to directly evolve the communication topology temporally.

Additionally, to perceive the system state and provide self-feedback ability, we devise mechanisms driven by both entropy and experience, then internalize them through preference optimization. First, we supervise the system state with the help of \textit{Entropy Distribution} to evaluate the uncertainty of the MAS, and we internalize it into the scheduler via reinforcement learning. Second, considering the self-feedback ability, we draw inspiration from \textit{Experience Learning}~\citep{zhang2025agent,tang2025chemagent}, and construct a retrieve-augment module to accumulate successful scheduling trajectories as experience. During training, we also internalize the experience, which retrieves corresponding historical scheduling trajectories as reference, and regularizes the current scheduling scheme.

Our Contributions can be summarized as:

\begin{itemize}[leftmargin=2em,itemsep=-0.1em]
\item[\ding{182}] We introduce \ourmethod, pioneering the exploration of Spatio-Temporal Evolving MAS. It can support dialogue-wise real-time scheduling through generative Spatio-Temporal evolution of communication topologies.
\item[\ding{183}] We devise a compact generative scheduler driven by Flow-Matching. It can learn underlying scheduling principles from entropy distribution and historical experience, thus possessing strong query-aware and dialogue-wise scheduling capability. 
\item[\ding{184}] Comprehensive experiments on 9 benchmarks from multiple real-world tasks demonstrate the state-of-the-art performance of our proposed \ourmethod. 
\end{itemize}
